\documentclass[12pt,a4paper]{article}
\usepackage[utf8]{inputenc}
\usepackage[T1]{fontenc}
\usepackage{amsmath,amssymb,amsthm}
\usepackage{geometry}
\usepackage{hyperref}
\usepackage{listings}
\usepackage{xcolor}
\usepackage{booktabs}
\usepackage{enumitem}
\usepackage{mdframed}
\usepackage{mathtools}
\usepackage{microtype}

\geometry{margin=2.5cm}

\lstset{
  basicstyle=\small\ttfamily,
  breaklines=true,
  frame=single,
  language=Mathematica,
  keywordstyle=\color{blue},
  commentstyle=\color{gray},
  backgroundcolor=\color{gray!5}
}

\newtheoremstyle{critical}{}{}{\itshape}{}{\bfseries\color{red}}{.}{ }{}
\theoremstyle{critical}
\newtheorem{issue}{Issue}

\newtheoremstyle{remark}{}{}{}{}{\bfseries}{.}{ }{}
\theoremstyle{remark}
\newtheorem{rem}{Remark}
\newtheorem{claim}{Claim}
\newtheorem{theorem}{Theorem}

\newtheoremstyle{verified}{}{}{}{}{\bfseries\color{blue!70!black}}{.}{ }{}
\theoremstyle{verified}
\newtheorem{verified}{Verified}

\title{\textbf{Critical Audit of SZ-DBC (Second Report):\\Detailed Balance Checker for Lattice MCMC}\\[1ex]
\large Correctness, rigor, and portability analysis\\
\normalsize Incorporating all fixes from the first report}
\author{Independent Technical Review}
\date{June 2026}

\begin{document}
\maketitle
\tableofcontents
\newpage

%===============================================================
\section{Scope and Method}
%===============================================================

This report constitutes a deep, independent audit of the SZ-DBC codebase as it
stands after the two rounds of fixes described in the previous report
(\texttt{critical\_report.tex}).  All seven issues raised there have been addressed;
this review therefore starts fresh and examines the \emph{current} code for
new correctness problems, rigour gaps, and undocumented limitations.

The methodology was: read every line of \texttt{dbc\_core.wl} and
\texttt{check.wls}; run all provided examples and verify their outputs; run
targeted Wolfram~Language micro-experiments to probe specific claims; and
reason from first principles about every non-trivial assertion made by the
code or its documentation.  All example runs were performed on the current
codebase with \texttt{wolframscript}.

\begin{mdframed}[backgroundcolor=green!8,linecolor=green!60!black,linewidth=1pt]
\textbf{Top-line verdict.}
The checker is now substantially more rigorous than the version audited in the
first report.  The critical RandomChoice weight bug is fixed, the D4 check is
algebraically exact, the ergodicity BFS starts from the correct seed, and
timed-out BFS paths abort with an error.  The EBE-based DB check is exact for
all algorithms with fewer than \texttt{ebeMaxK} branch conditions.

Two moderate correctness concerns survive: (1) an unverified assumption that
the post-D4 BFS introduces no new Piecewise conditions beyond those seen in the
pre-D4 BFS, silently relied upon when reusing EBE chambers; (2) the
\texttt{\$dbcFS} sentinel is treated as a DB violation rather than an analysis
failure.  Neither causes problems for the provided examples.  Three minor issues
round out the analysis.
\end{mdframed}

%===============================================================
\section{The BFS Engine: Correctness Review}
\label{sec:bfs}
%===============================================================

\subsection{Weight completeness}

A necessary condition for BFS correctness is that the leaf weights for any
single starting state sum to~1 (probability is conserved across all execution
paths).  This was verified experimentally on the first orbit representative of
\texttt{single\_metropolis.wl}:

\begin{lstlisting}
Total[#[[3]] & /@ leaves]  (* symbolic sum over all leaves *)
(* After Simplify: 1 *)
\end{lstlisting}

The sum is a large symbolic expression involving Piecewise coupling comparisons
and Boltzmann factors; \texttt{Simplify} reduces it to~$1$ exactly.  This
confirms that the accept/reject branches pair correctly and that the
weight-correction factors $2^k/n$ for \texttt{RandomChoice} are applied
consistently.  The proof that weight sums to~1 holds generally:

\begin{theorem}[Weight completeness]
For any algorithm whose random calls are all handled by the BFS engine (no
\texttt{cantHandle}), the leaf weights from any starting state sum to~$1$.
\end{theorem}

\begin{proof}[Sketch]
By construction, each \texttt{RandomReal[]} comparison branches into two paths
with weights $p$ and $1-p$ summing to~$1$.  Each \texttt{RandomChoice[list]}
with $n$ elements reads $k = \lceil\log_2 n\rceil$ bits; in-range leaves each
have weight $2^k/n$ (after correction), and the $2^k - n$ out-of-range paths
are discarded.  The $n$ in-range leaves carry weight $n \cdot (1/2^k) \cdot
(2^k/n) = 1$.  The \texttt{seqBernoulli} decomposition is a perfect partition
of~$[0,1]$ by construction (verified above).  The product of independent
branches gives a probability measure on the leaf set.
\end{proof}

\subsection{RandomChoice correction for non-power-of-2 pool sizes}

This was the critical bug fixed in the first revision.  The current code applies:
\begin{lstlisting}
weight *= 2^k/n;  (* after selecting valid index *)
\end{lstlisting}
for \texttt{RandomChoice[list]}, \texttt{RandomInteger[\{lo,hi\}]}, and the
vectorised forms.  For $n$ a power of~2 the factor is~$1$ and has no effect;
for general $n$ it corrects the per-leaf weight from $1/2^k$ to $1/n$.  The
broken examples \texttt{broken\_variable\_pool.wl} (pool size 3 or 4, $k=2$)
and \texttt{broken\_8way\_hop.wl} (pool size 7 or 8, $k=3$) now correctly
report \textbf{FAIL} with exact rational transition probabilities printed.

\begin{verified}
All nine examples produce the expected verdict.  D4 is no longer declared in
\texttt{\$symmetryGroup} by default (see Section~\ref{sec:d4perf}): the three
passing examples (\texttt{single\_metropolis}, \texttt{kawasaki},
\texttt{vmmc\_2d}) show \texttt{D4 rotational: not declared} in their summary
output.  DB correctness is unaffected.
\end{verified}

\subsection{Edge cases in the BFS engine}

\noindent\textbf{Empty range \texttt{RandomInteger[\{hi,lo\}]} with $hi < lo$.}
When $n = hi - lo + 1 \leq 0$, the code computes
$k = \texttt{IntegerLength}[n-1, 2]$.  For $n = 0$: \texttt{IntegerLength}$[-1,2]$
$= 1$, so one bit is read and $\texttt{val} \in \{0,1\}$; both satisfy
$\texttt{val} \geq n = 0$, so \texttt{\$dbc\$outOfRange} is thrown and the path
is silently discarded.  This is correct behaviour (a call with inverted bounds
has no valid outcomes) but is not documented.

\noindent\textbf{Single-element lists.}
\texttt{RandomChoice[\{x\}]} (n=1) returns $x$ immediately without reading bits
(explicit \texttt{n==1} branch).  \texttt{RandomInteger[\{lo,lo\}]} (n=1)
similarly returns \texttt{lo} with weight~$1$.  \texttt{seqBernoulli[\ldots,
\{e\}]} (one element) is short-circuited before the call site.  All correct.

\noindent\textbf{Zero-weight seqBernoulli guard.}
The guard \texttt{If[TrueQ[remainW == 0], Throw[\$dbc\$cantHandle[...]]]} was
added in the second revision.  Experimentally verified: for concrete zeros
(\texttt{ws = \{0,0,0\}}) and for symbolically simplified zeros
(\texttt{\{couplingJ[1,2,1], -couplingJ[1,2,1]\}}), Mathematica simplifies
\texttt{Total[ws]} to \texttt{0} and \texttt{TrueQ[0 == 0]} returns
\texttt{True}, so the guard fires.  For genuinely pathological cases where a
symbolic total equals zero but Mathematica does not simplify it, division by
zero would produce \texttt{Indeterminate}.  This is unlikely for physical
algorithms (weights are non-negative energy-difference expressions).

%===============================================================
\section{The EBE Exact DB Check}
\label{sec:ebe}
%===============================================================

\subsection{Phase 1: Condition extraction}

The extraction pattern matches \texttt{HoldPattern[Piecewise[cl\_List, \_]]} to
avoid premature evaluation of Piecewise clauses.  It retains only
\texttt{\_Less | \_LessEqual | \_Greater | \_GreaterEqual} conditions that
involve at least one symbolic parameter.  This correctly excludes \texttt{True},
\texttt{False}, \texttt{=== Infinity} sentinels (SameQ, not a comparison), and
numeric constants.

For \texttt{vmmc\_2d.wl}, link probabilities include clauses of the form
\texttt{\{1, eFwd === Infinity\}} using \texttt{===} (structural equality), which
are not matched by the comparison filter and are therefore never extracted as EBE
conditions.  Since \texttt{couplingJ[\ldots] === Infinity} evaluates to
\texttt{False} for any finite coupling, these clauses always contribute weight~0
and the extraction is correct.

\subsection{Phase 2: Chamber enumeration}

\begin{claim}[Correctness of arrangement BFS]
The chamber-adjacency BFS in Phase~2 visits every non-empty open chamber of the
hyperplane arrangement defined by the $k$ Piecewise conditions.
\end{claim}

\begin{proof}
The arrangement is finite; each hyperplane is defined by a linear inequality in
the coupling constants.  By the standard connectivity theorem (any two chambers
are path-connected via a sequence of codimension-1 face crossings), the
face-adjacency graph is connected.  The BFS therefore visits every chamber
reachable from the starting point.  By connectivity, this is all chambers.

For each adjacent chamber (sign-pattern obtained by flipping one condition),
\texttt{FindInstance[constraints, allSymParams, Rationals, 1]} is called.  For
rational-coefficient linear inequalities over the reals, a rational interior
point exists whenever the region is non-empty (rational points are dense in any
open set).  Thus if a chamber is non-empty, \texttt{FindInstance} will find a
representative.
\end{proof}

\noindent\textbf{Boundary omission.}
The open-chamber enumeration does not test the lower-dimensional faces where two
or more conditions are simultaneously tight (e.g., $J_1 = J_2$ exactly).  A DB
violation existing only on such a boundary (and not in any adjacent open chamber)
would be missed.  For Metropolis-type algorithms, transition probabilities are
continuous functions of the coupling constants; any violation in a positive-measure
region must therefore appear in an adjacent open chamber.  The only pathological
case is an algorithm with an explicit equality-trigger rule (\emph{e.g.}
\texttt{If[$J_1$ == $J_2$, ...]}), which would be non-standard.  This should be
documented.

\subsection{Phase 3: The \texorpdfstring{$\$dbcIsExpZero$}{dbcIsExpZero} check}
\label{sec:expzero}

After substituting a rational representative $J^\ast$ into every leaf weight,
the DB expression takes the form $\sum_i c_i e^{-\beta E_i}$ where $c_i$ are
rational and $E_i$ are rational multiples of $\beta$, \emph{i.e.}
$\sum_i c_i E^{-\beta r_i}$ with $r_i \in \mathbb{Q}$.  The function
\texttt{\$dbcIsExpZero} checks this is zero as follows:

\begin{enumerate}
\item \texttt{\$dbcMergeExp}: iteratively replaces
  $E^a \cdot E^b \to E^{a+b}$ until no such pairs remain
  (at most 30 iterations).
\item \texttt{\$dbcSplitTerm}: for each term, extract the unique Exp factor (if
  any) to produce \texttt{\{coefficient, exponent\}}.  Returns \texttt{\$dbcFS}
  if two or more distinct Exp factors remain after merging.
\item \texttt{\$dbcGroupByExp}: group terms by exponent using
  \texttt{Expand[\ldots] === 0}.
\item Check that the sum of rational coefficients in each group is zero.
\end{enumerate}

\noindent\textbf{Correctness for standard algorithms.}
Experimentally verified: after $J^\ast$ substitution, all Exp factors in
Metropolis-type algorithms are of the form $E^{-\beta r}$; \texttt{\$dbcMergeExp}
reduces any product $E^{-\beta r_1} \cdot E^{-\beta r_2}$ in one iteration.
The \texttt{\$dbcFS} branch is never triggered for any of the provided examples.
For algorithms producing genuine transcendental coefficients
(e.g.\ $E^{\sin(\beta)}$), merging would fail after 30 iterations and
\texttt{\$dbcFS} would be returned.

\begin{issue}[\texttt{\$dbcFS} treated as a violation, not an analysis failure]
\label{issue:dbcfs}
The check \texttt{If[res =!= True, violations[pair] = ...]} treats both
\texttt{False} (genuine DB violation) and \texttt{\$dbcFS} (analysis failure)
as violations.  A user whose algorithm produces a \texttt{\$dbcFS} result (e.g.\
an algorithm using $\sin(\beta)$ in the acceptance probability) would receive a
FAIL verdict with no diagnostic distinguishing the two cases.  The violation
record would contain a \texttt{tij} or \texttt{tji} field with an unresolved
symbolic expression, which is the only hint that analysis failed rather than DB
being violated.

This is a pre-existing issue from the first report, not yet fixed.  It is
harmless for all standard physics algorithms tested, but a defensive fix (checking
\texttt{res === \$dbcFS} separately and aborting with an error) is
straightforward.
\end{issue}

%===============================================================
\section{The D4 Algebraic Check: Review of the New Implementation}
\label{sec:d4}
%===============================================================

\subsection{Generator sufficiency}

The D4 check tests only two generators: \texttt{rotIdx=1} (90\(^\circ\) clockwise
rotation) and \texttt{rotIdx=4} (left-right reflection).  The claim is that
these generate the full D4 group.  This is verified:

\begin{proof}
With $r$ = rot90CW (order 4) and $s$ = reflLR (order 2), we have:
\begin{align*}
  r^0 &= \text{id},\quad r^1 = \text{rot90},\quad r^2 = \text{rot180},\quad r^3 = \text{rot270}\\
  s\cdot r^0 &= \text{reflLR},\quad s\cdot r^1 = \text{reflDiag},\quad
  s\cdot r^2 = \text{reflUD},\quad s\cdot r^3 = \text{reflAnti}
\end{align*}
These 8 elements are all distinct and cover D4.  Verified numerically by applying
all compositions to an asymmetric test state and confirming 8 distinct images
matching all \texttt{rotIdx} 0--7.
\end{proof}

The group-theoretic argument that verifying invariance for generators $r$ and $s$
implies invariance for all products is correct by induction: if $T(x\to y) =
T(r\cdot x\to r\cdot y)$ for all pairs, then applying this to the pair
$(s\cdot x, s\cdot y)$ gives $T(s\cdot x\to s\cdot y) = T(r\cdot s\cdot x\to
r\cdot s\cdot y)$, establishing invariance under $r\cdot s$.  By induction over
all words in $\{r, s\}$, invariance holds for all D4 elements.

\subsection{Correctness of Phase 3: checking $T(i\to j)$ vs $T(R\cdot i\to R\cdot j)$}

In Phase~3 of \texttt{\$dbcEBECheckD4}, the transition $T(R\cdot i\to R\cdot j)$
is looked up via \texttt{pairToLeafSrc[\{rsi, rsj\}]}, where \texttt{rsi} and
\texttt{rsj} are the indices of $R\cdot i$ and $R\cdot j$ in \texttt{allStates}.
The orbit expansion from the translation-only BFS covers all pairs in
\texttt{allStates} (since the translation group generates all $n^2$ positions).
Therefore $T(R\cdot i\to R\cdot j)$ is correctly assembled from the expanded
leaves.  If $T(R\cdot i\to R\cdot j) = 0$ and $T(i\to j) \neq 0$, the check
\texttt{\$dbcIsExpZero[tij - 0]} correctly detects the violation.

The check in Phase~3 evaluates the difference $T(i\to j) - T(R\cdot i\to
R\cdot j)$ using \texttt{\$dbcIsExpZero}.  Since both T values are sums of
rational $\times$ $E^{-\beta r}$ terms (after $J^\ast$ substitution), the
difference has the same structure.  The checker handles this correctly.

\subsection{Chamber reuse: a correctness assumption}
\label{sec:chamberreuse}

\begin{issue}[Chamber reuse assumption not verified]
\label{issue:chambers}
After D4 passes, the pipeline recomputes orbit representatives using the full
D4+translation group (reducing from 56 to 8 reps for \texttt{single\_metropolis})
and re-runs the BFS from these 8 reps.  The EBE DB check then reuses the 132
chambers computed during the D4 EBE check (which used the 56-rep BFS), skipping
Phase~2.

The correctness of this reuse relies on the following assumption:

\begin{mdframed}[backgroundcolor=yellow!20,linewidth=0.5pt]
\textbf{Reuse assumption.}  Every Piecewise condition appearing in the 8-rep BFS
leaves is also a condition in the 56-rep BFS leaves.  Equivalently, the 8-rep
hyperplane arrangement is a \emph{coarsening} of (or identical to) the 56-rep
arrangement.
\end{mdframed}

If this assumption fails, some 8-rep Piecewise conditions are not controlled by
the chamber decomposition; the $J^\ast$ for each chamber fixes them to specific
values, but those values may not be representative of the entire chamber.  The DB
check would be valid only at the single point $J^\ast$, not throughout the
chamber.

\textbf{Mathematical justification.}
If D4 invariance holds exactly (as verified by the algebraic D4 EBE check), then
any energy-difference comparison arising in the 8-rep BFS also arises in the
56-rep BFS (or is D4-equivalent to one that does, and since D4 acts on positions
and not types, the coupling atoms are preserved).  The 8-rep conditions are
therefore a subset of the 56-rep conditions, making the reuse valid.

\textbf{Empirical verification.}
In all provided examples, the EBE Phase~1 applied to the 8-rep BFS yields the
same $k$ as Phase~1 of the D4 EBE check:
\begin{itemize}
\item \texttt{single\_metropolis.wl}: $k = 18$ for both checks, 132 chambers
\item \texttt{vmmc\_2d.wl}: $k = 12$ for both checks, 512 chambers
\item \texttt{kawasaki.wl}: $k = 6$ for both checks, 12 chambers
\end{itemize}

\textbf{What is missing.}
The code does not assert that the 8-rep conditions form a subset of the 56-rep
conditions.  A defensive check would be:

\begin{lstlisting}
(* After Phase 1 of DB EBE, if reusing precomputedChambers: *)
If[!SubsetQ[allConds, precomputedConditions],
  Print["WARNING: new conditions not covered by precomputed chambers"];
  (* fall back to Phase 2 *)]
\end{lstlisting}

Without this check, a future algorithm where D4 invariance is algebraically
verified but the 8-rep BFS introduces new Piecewise conditions (e.g., from a
distance function the 56-rep never reached) would silently produce an incomplete
DB check.  The mathematical argument above rules this out for standard physics
algorithms, but a runtime assertion is prudent for a tool claiming exactness.
\end{issue}

\subsection{C4-only case}

The same reuse concern applies when only C4 (rotations, not reflections) is
verified.  In this case, the post-C4 BFS uses C4+translation orbits (between 8
and 56 reps, depending on the system).  The D4 EBE check computed chambers from
the 56-rep BFS.  For C4-symmetric algorithms without full D4, it is possible
(though unlikely for standard Hamiltonians) that the post-C4 BFS introduces
conditions absent from the 56-rep BFS.

%===============================================================
\section{The Translational Invariance Check}
\label{sec:tau}
%===============================================================

\subsection{Correctness of weight-only check}

The $\tau$-BFS verifies that leaf weights are free of $\tau_r$ and $\tau_c$.  It
does not verify that the next-state for each leaf is the correct $\tau$-shifted
version of the non-augmented next state.  For algorithms where spatial decisions
use only pairwise differences, the state update (e.g., moving a particle from
$p$ to $p + \Delta$) always produces states that correctly carry $\tau$.  A
pathological algorithm could, in principle, produce $\tau$-free weights while
constructing next states that are not translates of the rep's next states.  Such
an algorithm would pass the $\tau$-check while being non-translation-invariant.
No such algorithm is physically meaningful, but the limitation should be noted.

\subsection{The cantHandle path for quadratic\_field.wl}

Confirmed from the current run: \texttt{quadratic\_field.wl} triggers an error
\begin{lstlisting}
RandomInteger[{{-16,16},4}]: 24 bits needed --- use count x RandomInteger[{lo,hi}]
\end{lstlisting}
and falls back to the non-translation-invariant path.  This is the improved error
message from the second revision.  The actual mechanism is that Mathematica's
symbolic evaluation of \texttt{Mod[pos, n]} with a $\tau$-augmented position
internally calls \texttt{RandomInteger[\{-16,16\}, 4]}, which is intercepted.
The checker correctly falls back and the DB check is exact.

\subsection{Soundness of checking from orbit representatives}

\begin{theorem}[Sufficiency of orbit representatives for $\tau$-check]
If the $\tau$-BFS produces $\tau$-free leaf weights for each translation-orbit
representative, it produces $\tau$-free weights for every state.
\end{theorem}

\begin{proof}
Every state $s$ is a translation $T_{(dr,dc)}(\text{rep})$ for some
$(dr,dc) \in \{0,\ldots,n-1\}^2$.  Running the algorithm from $s + \tau =
T_{(dr,dc)}(\text{rep}) + \tau = T_{(dr,dc)}(\text{rep} + \tau)$ (where
$\tau = (\tau_r, \tau_c)$) is equivalent to running from $\text{rep} + \tau'$
with $\tau' = \tau + (dr, dc)$.  Since the $\tau$-BFS checks rep at offset $\tau$
(a free variable), the output subsumes all integer offsets of $\tau$.  Hence
$\tau$-freedom for the rep implies $\tau$-freedom for all states in its orbit.
\end{proof}

%===============================================================
\section{Ergodicity Check}
\label{sec:ergodicity}
%===============================================================

The ergodicity BFS was fixed in the first revision to start from the correct
\texttt{seedState} index.  The current implementation looks up
\texttt{seedIdx = Lookup[stateToIdx, Key[seedState], 1]} and falls back to
index~1 if the seed is not found.  This fallback cannot occur for well-formed
inputs (the seed state is always in \texttt{allStates}) but provides a safe
default.

\textbf{Interaction with maxDepth.}
The BFS leaf set may be incomplete if any transition requires more than
\texttt{maxDepth}~bits (currently 22).  An algorithm with a long-tailed random
path (e.g., a geometric distribution with small success probability) might have
its low-probability transitions truncated.  In the current code, reaching
\texttt{maxDepth} causes a hard error (exit) rather than a silent omission, so
the ergodicity check cannot silently produce false PASSes from incomplete leaves.

%===============================================================
\section{New and Undocumented Issues}
\label{sec:newissues}
%===============================================================

\subsection{NormalDistribution discretisation: undocumented truncation error}
\label{sec:normaldist}

The \texttt{RandomVariate[NormalDistribution[$\mu$, $\sigma$]]} handler
discretises onto $\{\texttt{Round}(\mu)-\texttt{nMax}, \ldots,
\texttt{Round}(\mu)+\texttt{nMax}\}$ where
\texttt{nMax}~$=$~\texttt{Floor}[\texttt{\$dbcCurrentNGrid}/2].  For a
$3\times3$ grid this gives $\texttt{nMax}=1$ and only 3 values: $\{\mu-1, \mu,
\mu+1\}$.  The probability mass outside this range is discarded, and the weights
are renormalised.

This introduces a systematic error:
\begin{center}
\begin{tabular}{lll}
\toprule
$\sigma$ & Mass covered & Mass discarded \\
\midrule
0.5 & 99.7\% & 0.3\% \\
1.0 & 86.6\% & 13.4\% \\
2.0 & 54.7\% & \textbf{45.3\%} \\
\bottomrule
\end{tabular}
\end{center}
The renormalisation ensures the BFS leaf weights still sum to~1, but the
transition probabilities assigned to each outcome are \emph{wrong}: the checker
verifies DB for a truncated distribution, not for the actual NormalDistribution.
A user applying the checker to an algorithm that genuinely uses
\texttt{RandomVariate[NormalDistribution[$\mu$, $\sigma$]]} with $\sigma \geq 1$
on a $3\times3$ grid would receive a DB verdict that is \textbf{not valid for the
intended distribution}.

The correct fix is either to use the exact CDF-based weights (which is what the
code already does for the truncated range) \emph{and} to include a tail-capture
term that catches all values outside $[\mu-\texttt{nMax}, \mu+\texttt{nMax}]$,
or to document explicitly that the checker is approximating the distribution and
that $\sigma \ll \texttt{nMax}$ is required for the approximation to be
acceptable.  The current code prints no warning, making this a silent
correctness risk.

\subsection{Coupling atom arity constraint}
\label{sec:arity}

The validation check
\begin{lstlisting}
!AllTrue[symCouplings, MatchQ[#, _[_Integer, _Integer, _Integer]]&]
\end{lstlisting}
requires every coupling atom to have \emph{exactly} three integer arguments.  An
algorithm using two-argument couplings (\texttt{coupling[a,b]}, for
distance-independent interactions) would fail with the message ``coupling atom
arguments must all be integers,'' which does not explain the three-argument
requirement.  This is an unnecessary constraint: the coupling symmetry rule
\texttt{h[b\_, a\_, rest\_\_Integer] /; b > a := h[a, b, rest]} requires
\texttt{rest} to match at least one argument, but a two-argument coupling could
be supported with a separate rule.  The error message should at minimum clarify
that three integer arguments are required, or the constraint should be relaxed.

\subsection{Legacy dead code}
\label{sec:dead}

The functions \texttt{\$dbcApplyRot90} and \texttt{\$dbcApplyReflect} are defined
in \texttt{dbc\_core.wl} but never called in the current pipeline (the D4 action
is computed everywhere via \texttt{\$dbcApplyGElem} and \texttt{\$dbcD4Pos}).
These are artefacts of an earlier version and constitute dead code.  Their
presence is harmless but confusing; they should be removed.

\subsection{RandomInteger with inverted bounds}
\label{sec:inverted}

As noted in Section~\ref{sec:bfs}, \texttt{RandomInteger[\{hi, lo\}]} with
$hi < lo$ produces $n \leq 0$.  The code silently discards all paths via
\texttt{outOfRange} rather than raising a \texttt{cantHandle} error.  A user
whose algorithm contains an inverted-bound call due to a bug would not be informed
by the checker.  A guard
\begin{lstlisting}
If[n <= 0, Throw[$dbc$cantHandle["RandomInteger: invalid range"], $dbc$tag]]
\end{lstlisting}
would make the failure visible.

%===============================================================
\section{Correctness of Orbit Expansion}
\label{sec:orbits}
%===============================================================

\subsection{The orbit map construction}

In \texttt{\$dbcComputeOrbits}, the canonical representative \texttt{minImg} is
chosen as the lexicographically smallest state in the orbit.  The orbit map
\texttt{orbitMap[s] = g} satisfies $g(\text{minImg}) = s$.  The comment in the
source correctly identifies the bug that was fixed: when the initially-visited
state $s \neq \texttt{minImg}$, the orbit map must be computed \emph{from
minImg}, not from $s$.  The fix (recomputing \texttt{imgPairsFromRep} from
\texttt{minImg}) is in place and correct.

\subsection{Orbit expansion correctness}

The expansion $T(\texttt{sKey} \to g(\texttt{nextState})) = T(\texttt{minImg} \to \texttt{nextState})$ relies on
the G-invariance of the algorithm, which is established independently by the
$\tau$-check and D4 EBE check before the expansion is used.  All states in
\texttt{allStates} are correctly mapped under D4 actions (verified: zero misses
for all 504 states and all 8 D4 elements for \texttt{single\_metropolis}).

\subsection{Self-loops and T-matrix diagonal}

Pairs where $i = j$ (self-loops, \emph{i.e.}\ the algorithm proposes a rejected
or null move and returns the same state) are excluded from the DB check by
\texttt{pairsArr = Select[pairsArr, \#[[1]] =/= \#[[2]] \&]}.  The DB condition
$T(i\to i) \cdot \pi(i) = T(i\to i) \cdot \pi(i)$ is trivially satisfied, so
this exclusion is correct.

%===============================================================
\section{Summary of Open Issues}
\label{sec:summary}
%===============================================================

\begin{table}[h!]
\centering
\begin{tabular}{lllll}
\toprule
\textbf{Issue} & \textbf{Severity} & \textbf{Status} & \textbf{Affects examples?} & \textbf{Section} \\
\midrule
\$dbcFS treated as violation & Moderate & Open & No & \ref{sec:expzero} \\
Chamber reuse assumption & Moderate & Open (undocumented) & No & \ref{sec:chamberreuse} \\
NormalDistribution truncation & Significant & Open (undocumented) & No & \ref{sec:normaldist} \\
Coupling arity constraint & Minor & Open & No & \ref{sec:arity} \\
Inverted RandomInteger bounds & Minor & Open & No & \ref{sec:inverted} \\
Dead code (ApplyRot90 etc.) & Cosmetic & Open & No & \ref{sec:dead} \\
D4 default performance & Performance & \textbf{Fixed} & Yes (vmmc\_2d: 332s$\to$217s) & \ref{sec:d4perf} \\
Probabilistic SZ fallback & Correctness & \textbf{Fixed} & No & \ref{sec:d4perf} \\
\bottomrule
\end{tabular}
\caption{Issues after both previous revision rounds, plus performance findings.
  ``Affects examples?'' refers to the nine provided example algorithms.
  The bottom two rows describe changes made in the current revision.}
\end{table}

The highest-severity open issue is the NormalDistribution truncation, which
produces \textbf{wrong transition probabilities without warning} for
$\sigma \geq 1$ on small grids.  Algorithms using Gaussian displacement proposals
should not be verified with this checker without explicitly confirming that $\sigma
\ll \texttt{Floor}[\texttt{nGrid}/2]$.

%===============================================================
\section{Performance Analysis: The D4 Orbit-Reduction Trade-off}
\label{sec:d4perf}
%===============================================================

\subsection{Background and motivation}

The original rationale for including \texttt{"D4"} in \texttt{\$symmetryGroup}
was that, if D4 symmetry holds, orbit representatives can be reduced from 56
(translation only) to 8 (D4$+$translation).  Fewer representatives were expected
to reduce the number of state pairs in EBE Phase~3 by a factor of~7, making the
DB check 7$\times$ faster.  This section shows that this reasoning was
quantitatively incorrect for VMMC, and documents the resulting architectural
change.

\subsection{Timing breakdown with and without D4}

Comprehensive profiling of \texttt{vmmc\_2d.wl} ($n=3$, three particle types)
produced the following wall-clock breakdown.

\paragraph{Original behaviour (D4 included in \texttt{\$symmetryGroup}).}
\begin{center}
\begin{tabular}{lrl}
\toprule
\textbf{Step} & \textbf{Time} & \textbf{\% of total} \\
\midrule
$\tau$-BFS (56 reps) & 12.8s & 4\% \\
BFS (56 reps) & 11.7s & 4\% \\
D4 EBE check (Phase 2 + Phase 3 $\times$ 2 generators) & 203.3s & 61\% \\
BFS (8 D4-reduced reps) & 1.9s & $<$1\% \\
Ergodicity & 1.9s & $<$1\% \\
DB EBE check (Phase 3, 8-rep leaves) & 101.1s & 30\% \\
\midrule
\textbf{Total} & \textbf{332s} & \\
\bottomrule
\end{tabular}
\end{center}

\paragraph{New behaviour (D4 removed; translation only).}
\begin{center}
\begin{tabular}{lrl}
\toprule
\textbf{Step} & \textbf{Time} & \textbf{\% of total} \\
\midrule
$\tau$-BFS (56 reps) & 13.9s & 6\% \\
BFS (56 reps) & 12.8s & 6\% \\
D4 EBE check & 0s & --- \\
Ergodicity & 2.0s & 1\% \\
DB EBE check (Phase 3, 56-rep leaves) & 181.7s & 84\% \\
\midrule
\textbf{Total} & \textbf{217s} & \\
\bottomrule
\end{tabular}
\end{center}

The runtime decreases from 332s to 217s, a \textbf{1.53$\times$ speedup}, by
removing two lines from the algorithm file.  For \texttt{single\_metropolis.wl}
($k=18$, 132 chambers) the improvement is from $\approx$65s to $\approx$45s.

\subsection{Why D4 orbit reduction barely helps Phase 3}

The expected 7$\times$ speedup from D4 orbit reduction did not materialise because
the cost of EBE Phase~3 is dominated by the number of \emph{communicating state
pairs}, not the number of orbit representatives.

The pair array \texttt{pairsArr} is computed by expanding orbit-rep leaves through
the full orbit map and deduplicating by symmetry.  With 56 translation-only reps,
the translation group of size 9 is used for deduplication; with 8 D4+translation
reps, the group of size 72 is used.  The resulting pair counts are:
\begin{itemize}
  \item 56 translation reps $\to$ 5{,}688 distinct pairs (up to translation)
  \item 8 D4+translation reps $\to$ 5{,}688 distinct pairs (up to D4+translation)
\end{itemize}

The pair count is \emph{identical}.  Both representations cover all 504 states
and the same set of nonzero transitions; what changes is only the group used to
deduplicate.  Since the state space is fully ergodic and the transition graph is
dense, every pair reachable from any single rep is also reachable from any other,
so deduplication by a larger group yields the same universal set of pairs.

The only Phase~3 cost that differs between 56- and 8-rep runs is Step~A, the
$J^\ast$ substitution into leaf weights: 4{,}368 leaves (56-rep) versus 648 leaves
(8-rep).  Empirically this costs $\approx$164\,ms/region versus $\approx$17.5\,ms/region
--- a 9.4$\times$ increase for a 6.7$\times$ increase in leaf count.
Over 512 regions this is 74.7\,s extra, not the 7$\times$ saving on the pair loop
that was expected.

\subsection{Why the D4 check appeared cheap (a subtle asymmetry)}

The D4 EBE check also runs Phase~3 with 56-rep leaves, but costs only
$\approx$203s despite checking 2 generators instead of 1.  This is because its
per-pair work is structurally much lighter than the DB check.  In the D4 check,
the test expression is $T(i\to j) - T(Ri\to Rj)$.  For a D4-symmetric algorithm
these two T-values have identical symbolic structure after $J^\ast$ substitution,
so the \texttt{Expand} step reduces the expression to~0 immediately and
\texttt{\$dbcIsExpZero} is never called.  In the DB check, the expression
$T(i\to j)\cdot E^{-\beta E_i} - T(j\to i)\cdot E^{-\beta E_j}$ requires a full
rational-coefficient grouping by Boltzmann exponent on every non-trivial pair.
The D4 Phase~3 benefits from early termination in a way the DB Phase~3 cannot.

\subsection{Implication: D4 is not a performance optimisation for dense transition graphs}

The D4 orbit reduction was designed around a cost model that assumes Phase~3
scales with the number of orbit representatives.  That assumption holds when the
transition graph is sparse (each rep connects to a small, non-overlapping
neighbourhood).  For ergodic algorithms on a 3$\times$3 grid the transition graph
is dense and the assumption fails.

D4 verification remains available as an opt-in correctness check --- useful when
one specifically wants to verify rotational symmetry of the algorithm --- but
should \textbf{not} be included in \texttt{\$symmetryGroup} for performance
reasons.

\subsection{The probabilistic fallback: permanently disabled}

The \texttt{ebeMaxK} parameter previously defaulted to~50, which would trigger
a probabilistic Schwartz-Zippel fallback for any algorithm with more than 50
Piecewise conditions.  This default has been raised to~10{,}000 (effectively
unlimited) to ensure the checker always runs the exact EBE algorithm.  To use the
probabilistic path explicitly, pass \texttt{-mode SZPure} or lower
\texttt{-ebeMaxK} on the command line.

%===============================================================
\section{Python and Open-Source Portability Analysis}
\label{sec:portability}
%===============================================================

\subsection{Motivation and scope}

SZ-DBC is implemented in Wolfram Mathematica, which requires a proprietary
licence.  This severely limits adoption.  The previous \texttt{report.tex}
discussed language portability at a high level; the codebase has since changed
substantially (FullSimplify is no longer the primary DB check, having been
replaced by EBE).  This section provides a detailed, up-to-date analysis of what
can and cannot be ported to Python, Julia, or C++.

\subsection{The five core mechanisms and their portability}

\subsubsection{Mechanism 1: Dynamic function interception (\texttt{Block})}

\texttt{Block[\{RandomReal = f, RandomInteger = g, ...\}, alg[state]]} rebinds
every call to \texttt{RandomReal} and \texttt{RandomInteger} in the dynamic scope
of \texttt{alg[state]}, regardless of call depth.  The algorithm under test
requires no modification.

\textbf{Python.}  Python provides no equivalent at the standard language level.
\texttt{unittest.mock.patch} can intercept named function calls, but only for
module-level attributes, not for arbitrary call trees.  Crucially, if the
algorithm calls \texttt{import random; random.random()}, patching
\texttt{random.random} in the test harness does not intercept calls made through
a different import path.  Algorithms using NumPy (\texttt{np.random.rand()}) or
multiple RNG sources cannot be intercepted uniformly.  Any call that escapes the
interception returns a concrete float, and the BFS silently produces an incorrect
leaf weight with no warning.

The workaround --- requiring the algorithm to accept an explicit RNG object --- is
the only viable Python strategy, but it is \emph{invasive}: the algorithm must be
rewritten to pass \texttt{rng=symbolic\_rng} through all function calls.  This
changes the algorithm's interface, introduces a risk of translation errors, and
breaks the key property that ``the algorithm requires no modification.''

\textbf{Julia.}  Cassette.jl provides a true call-tree overdubbing mechanism.  A
\texttt{Cassette.@overdub ctx rand()} call can redirect every invocation of
\texttt{rand()} (and \texttt{rand(T)}, \texttt{rand(rng, ...)}) at the LLVM-IR
level, regardless of call depth or module.  This is semantically equivalent to
Mathematica's \texttt{Block}.  It is the single most important feature for a
rigorous port.  However, Cassette.jl is a complex, low-level package with a steep
learning curve and occasional compilation instability on newer Julia versions.

\textbf{C++.}  No standard mechanism exists.  One can compile the algorithm with a
special RNG header that is swapped at link time, but this requires build-system
cooperation and prevents use of third-party RNG implementations.

\subsubsection{Mechanism 2: Operator overloading for comparison tokens}

When \texttt{RandomReal[]} is intercepted, the returned symbolic token
\texttt{\$dbc\$irand[j, acceptTestI]} must participate in comparison expressions
like \texttt{token < Exp[-$\beta$ dE]}.  In Mathematica, this is handled via
\emph{UpValues}: pattern-matching rules attached to the token head that trigger
when the token appears in a \texttt{Less} expression.

\textbf{Python.}  This is portable.  A Python class \texttt{IntervalToken} can
define \texttt{\_\_lt\_\_(self, other)} to call the BFS's \texttt{accept\_test}
method, returning \texttt{True} or \texttt{False} after reading a bit and updating
the weight.  The algorithm's conditional
\texttt{if token < threshold:} triggers \texttt{\_\_lt\_\_} transparently.  This
works for direct comparisons; arithmetic (\texttt{token + c < threshold}) requires
additional \texttt{\_\_add\_\_} overloading.

\textbf{Julia.}  Also portable via method overloading:
\texttt{Base.:$<$(token::IntervalToken, p)} is syntactically analogous to the
Mathematica UpValue.

\textbf{C++.}  Operator overloading is native to C++.  The \texttt{operator<}
method on a custom \texttt{IntervalToken} type works transparently.

\subsubsection{Mechanism 3: Symbolic algebra for Piecewise expressions}

After the BFS, leaf weights are symbolic expressions involving
\texttt{Piecewise} conditions on coupling constants and Boltzmann factors
$E^{-\beta E_i}$.  The EBE Phase~1 extracts these conditions, Phase~2 enumerates
chambers, and Phase~3 substitutes rational representatives and evaluates the
\texttt{\$dbcIsExpZero} check.

\paragraph{\$dbcIsExpZero in Python.}  This specific check (\emph{group by
$E^{-\beta r}$ exponent, sum rational coefficients, verify zero}) requires only
rational arithmetic and dictionary operations.  After substituting rational $J^\ast$
values, all $E_i$ are rational and all $c_i$ are rational.  This is fully
implementable in Python using \texttt{fractions.Fraction} and a dict keyed by
the rational exponent.  \textbf{This is the easiest component to port.}

\paragraph{Piecewise handling.}  During the BFS (before $J^\ast$ substitution),
leaf weights contain nested \texttt{Piecewise} expressions.  Mathematica's
\texttt{PiecewiseExpand} automatically computes the Cartesian product of all
Piecewise conditions in a product expression.  SymPy's
\texttt{piecewise\_fold} does not perform this cross-product expansion; it only
flattens nested Piecewise into the same Piecewise, not products of two
independent Piecewise.  A Python port must implement cross-product expansion
manually, which is a non-trivial algebraic operation.

\paragraph{Condition feasibility.}  In Phase~2, Mathematica's
\texttt{FindInstance[constraints, vars, Rationals, 1]} is a complete decision
procedure for systems of linear rational inequalities, returning an exact rational
point.  This has no Python equivalent.  SymPy's \texttt{Poly.all\_roots} and
\texttt{solve\_univariate\_inequality} handle some cases but are not complete for
multivariate systems.  The most practical replacement:
\begin{enumerate}
\item Use \texttt{scipy.optimize.linprog} to find a floating-point interior point.
\item Convert to a rational approximation using
  \texttt{fractions.Fraction.limit\_denominator(N)} for large $N$.
\item Verify the rational point satisfies all strict inequalities.
\end{enumerate}
This is an engineering solution that works in practice but can fail in degenerate
cases (e.g., near the boundary of a very narrow chamber).  Mathematica's
\texttt{FindInstance} gives an \emph{exact} rational point by Fourier-Motzkin
elimination.

\paragraph{Symbolic expression maintenance.}  During the BFS, weights are
products of acceptance probabilities, which involve Piecewise expressions in the
coupling constants.  Mathematica's pattern-matching evaluator automatically
simplifies these products; SymPy requires explicit calls to \texttt{simplify} or
\texttt{cancel} after each multiplication.  For complex algorithms (e.g., VMMC),
the symbolic expressions can grow very large; Mathematica's evaluator handles
these via efficient internal normal forms that SymPy cannot match.

\subsubsection{Mechanism 4: State enumeration and orbit computation}

These are combinatorial algorithms with no symbolic algebra.

\textbf{Python, Julia, C++.}  All fully portable.  \texttt{itertools.permutations}
(Python), comprehensions (Julia), or STL containers (C++) handle state
enumeration.  The orbit computation is graph-theoretic and translates directly.

\subsubsection{Mechanism 5: EBE Phase 2 chamber enumeration}

The chamber-adjacency BFS is a graph traversal algorithm.  The only Mathematica
dependency is \texttt{FindInstance} (discussed above).  The BFS logic itself is
trivially portable.

\subsection{Component-by-component portability summary}

\begin{table}[h!]
\centering
\begin{tabular}{lllll}
\toprule
\textbf{Component} & \textbf{Python} & \textbf{Julia} & \textbf{C++} & \textbf{Notes} \\
\midrule
Block interception & \textbf{No} & Via Cassette.jl & Template mock & Key limitation \\
Operator UpValues & Yes & Yes & Yes & Straightforward \\
\$dbcIsExpZero & Yes & Yes & Yes & Pure rational arith. \\
PiecewiseExpand & Partial & Partial & No & Manual cross-product \\
FindInstance & Approx. & Approx. & Via LP solvers & Not exact \\
State enumeration & Yes & Yes & Yes & Trivial \\
Orbit computation & Yes & Yes & Yes & Trivial \\
BFS engine (queue) & Yes & Yes & Yes & Trivial \\
Coupling symmetry rule & Yes & Yes & Yes & Trivial \\
seqBernoulli & Yes & Yes & Yes & Straightforward \\
\bottomrule
\end{tabular}
\caption{Portability of each SZ-DBC component.  ``Yes'' = direct port possible.
``Partial'' = feasible with additional implementation effort.
``Approx.'' = feasible with reduced exactness guarantees.
``No'' = fundamental limitation.}
\end{table}

\subsection{What the removal of FullSimplify changes}

The original code (as described in \texttt{report.tex}) used Mathematica's
\texttt{FullSimplify} as the primary DB check.  This was the component most
deeply tied to Mathematica's proprietary algebra engine: no other CAS matches
Mathematica's \texttt{FullSimplify} in breadth or reliability for mixed
exponential-Piecewise expressions.

The current code uses EBE as the primary check, with FullSimplify available only
as an optional mode (\texttt{-mode FullSimplify}).  The EBE check's inner
computation (\texttt{\$dbcIsExpZero}) is pure rational arithmetic --- no CAS
required.  \textbf{This is the most significant portability improvement.}  The
EBE strategy makes the ``heavy CAS'' component optional rather than mandatory.

However, the EBE approach introduces its own Mathematica dependency:
\texttt{FindInstance} for chamber enumeration.  Replacing this with an approximate
LP-based solution degrades the exactness guarantee.  The SZ fallback (probabilistic
check at random coupling points) is fully portable but sacrifices the completeness
guarantee.

\subsection{A realistic Python port: what would it look like?}

A realistic Python port of SZ-DBC would have the following structure:

\begin{enumerate}
\item \textbf{Algorithm interface:} The algorithm must be written (or rewritten)
  to accept an explicit RNG object.  A \texttt{SymbolicRNG} class provides
  \texttt{random\_real()} (returning an \texttt{IntervalToken}),
  \texttt{random\_choice(items)} (using seqBernoulli or rejection sampling),
  and \texttt{random\_integer(lo, hi)} (rejection sampling).

\item \textbf{BFS engine:} A Python queue manages bit-string prefixes.  The
  algorithm is called with each prefix; the \texttt{IntervalToken} class's
  operator overloads read bits and accumulate weights.

\item \textbf{Symbolic leaf weights:} After BFS, leaf weights are Python
  \texttt{Piecewise} objects (a list of \{condition, value\} pairs with a
  default).  The condition is a Python expression tree over the coupling atoms.

\item \textbf{EBE Phase 1:} Traverse all leaf weights and collect unique
  conditions (linear inequalities in coupling atoms) using symbolic expression
  equality.  SymPy can do this, albeit slowly for large expression trees.

\item \textbf{EBE Phase 2:} Chamber BFS using \texttt{scipy.optimize.linprog}
  for interior-point finding, with rational conversion.

\item \textbf{EBE Phase 3:} Substitute rational $J^\ast$ into leaf weights;
  evaluate Piecewise conditions; group by exponent; sum rational coefficients.
  This is pure Python with \texttt{fractions.Fraction}.
\end{enumerate}

\textbf{Limitations of this port:}
\begin{itemize}
\item The algorithm must be rewritten to use the explicit RNG interface.  Any
  algorithm that calls Python's \texttt{random} or NumPy's RNG directly would
  bypass the BFS.  This is both a usability burden and a correctness risk (missed
  random calls produce silently wrong transition probabilities).
\item EBE Phase~2 is approximate (LP-based rational approximation), not exact.
  Chamber boundaries in very narrow regions might be misclassified.
\item Piecewise cross-product expansion must be implemented manually.  SymPy's
  \texttt{piecewise\_fold} is insufficient.
\item The Python port would be significantly slower than the Mathematica version
  due to Python's interpreted overhead and SymPy's slower CAS.
\end{itemize}

\subsection{Julia: the recommended open-source alternative}

Julia with Cassette.jl + Symbolics.jl is the most viable open-source
implementation, for the following reasons:

\begin{enumerate}
\item \textbf{Cassette.jl overdubbing} provides a true Block-like interception
  mechanism for \texttt{rand()} at any call depth.  The algorithm requires zero
  modification.  This is the single most important feature.

\item \textbf{Multiple dispatch} in Julia is semantically equivalent to
  Mathematica's UpValue system for operator overloading.  Defining
  \texttt{Base.:$<$(token::IntervalToken, p)} is a direct translation.

\item \textbf{Symbolics.jl} provides symbolic algebra with significantly faster
  compile-time expression evaluation than SymPy.  \texttt{substitute} and
  \texttt{expand} handle the rational-coefficient Boltzmann algebra.

\item \textbf{\$dbcIsExpZero} translates directly using Julia's \texttt{Dict}
  and \texttt{Rational\{Int\}} types.

\item \textbf{FindInstance} remains a gap.  Julia's JuMP.jl can find LP interior
  points; converting to rationals requires care.  An exact alternative would be to
  implement Fourier-Motzkin elimination in Julia.
\end{enumerate}

\textbf{Remaining challenges in Julia:}
\begin{itemize}
\item Cassette.jl has had API instability across Julia versions; the overdubbing
  context for \texttt{rand()} must be carefully constructed to handle all forms
  (\texttt{rand()}, \texttt{rand(Float64)}, \texttt{rand(rng, ...)}).
\item Symbolics.jl has limited Piecewise support; cross-product expansion must be
  implemented from scratch.
\item The seqBernoulli decomposition and interval-tracking logic must be
  reimplemented; no direct equivalent exists in any Julia library.
\end{itemize}

\subsection{C++: not recommended}

A C++ port would require a complete implementation of symbolic expression trees,
Piecewise condition tracking, and rational arithmetic.  While feasible (using
e.g.\ the GiNaC computer algebra library), it would involve tens of thousands of
lines of code with no natural analogue of Mathematica's pattern-matching evaluator.
The effort is not justified given the availability of Julia as a better
alternative.

\subsection{The fundamental limitation: algorithm transparency}

The deepest limitation of any non-Mathematica port is that \textbf{the checker
cannot guarantee that all random calls are intercepted}.  In Mathematica, the
Block mechanism intercepts every evaluation of the bound symbol, regardless of
nesting.  In Python, any call that bypasses the explicit RNG interface (e.g., a
\texttt{random.random()} call inside a helper imported from a third-party library)
returns a concrete float and produces a silently incorrect transition weight.  In
Julia (with Cassette.jl), the overdubbing mechanism can intercept calls at the
compiled bytecode level, providing a much stronger guarantee --- but only if all
random calls go through \texttt{rand()} or \texttt{Random.rand()}, not through
e.g.\ a hand-written LCG in a C extension.

This issue is \textbf{not resolvable} by any finite amount of engineering: it is
a consequence of Python's (and C++'s) lack of dynamic rebinding at the language
semantics level.  In practice, it means that a Python port is inherently less
trustworthy than the Mathematica version for algorithms that use multiple RNG
sources or third-party libraries with their own random-number calls.

\textbf{Mitigation.}  The Python port should include a mechanism to detect
``escaped'' random calls.  One approach: wrap the algorithm in a context manager
that monkey-patches \texttt{random.random}, \texttt{random.randint}, and
\texttt{numpy.random.rand} (and their common aliases) at the module level, and
fails loudly if any of these is called without going through the interceptor.
This is fragile but provides a best-effort guarantee.

%===============================================================
\section{Recommendations}
\label{sec:recommendations}
%===============================================================

\subsection{Immediate fixes}

\begin{enumerate}
\item \textbf{Add NormalDistribution warning.}  Before using the truncated
  discretisation, print: ``WARNING: NormalDistribution discretised to $\pm$nMax
  values around $\mu$; $X$\% of probability mass discarded. Results are exact
  only for $\sigma \ll$ nMax.''

\item \textbf{Fix \$dbcFS false-positive.}  In the EBE Phase~3 check:
  \begin{lstlisting}
  Which[res === True, Null,
        res === $dbcFS, Return[$dbc$cantHandle["$dbcIsExpZero failed..."]],
        True, violations[pair] = ...]
  \end{lstlisting}

\item \textbf{Add chamber-reuse subset assertion.}  When reusing precomputed
  chambers, verify that the new BFS conditions form a subset:
  \begin{lstlisting}
  If[!SubsetQ[allConds, precomputedConds],
    Print["INTERNAL ERROR: new conditions not covered by chambers"]; Exit[1]]
  \end{lstlisting}

\item \textbf{Guard inverted RandomInteger bounds.}  Add:
  \begin{lstlisting}
  If[n <= 0, Throw[$dbc$cantHandle["RandomInteger: hi < lo"], $dbc$tag]]
  \end{lstlisting}

\item \textbf{Improve coupling arity message.}  Change the error message to
  explicitly state that three integer arguments are required.
\end{enumerate}

\subsection{Documentation}

\begin{enumerate}
\item Document the open-chamber omission (measure-zero boundaries not checked).
\item Document that the $\tau$-check verifies transition weights but not next-state
  correctness.
\item State explicitly that the NormalDistribution handler is an approximation.
\item Document the chamber-reuse assumption (D4 invariance implies condition
  subset, not verified at runtime).
\end{enumerate}

\subsection{For a Python or Julia port}

\begin{enumerate}
\item \textbf{Julia (recommended):} Use Cassette.jl for RNG interception,
  Symbolics.jl for expression algebra, and implement \$dbcIsExpZero and
  seqBernoulli from scratch.  Accept LP-based (approximate) chamber enumeration
  as an engineering trade-off.
\item \textbf{Python:} Require the algorithm to accept an explicit
  \texttt{symbolic\_rng} object; implement monkey-patching of common RNG
  module-level functions as a best-effort interception guard; use
  \texttt{fractions.Fraction} for exact coefficient arithmetic; and use
  \texttt{scipy.optimize.linprog} for chamber enumeration with rational
  conversion.  Document that the interception guarantee is weaker than the
  Mathematica version.
\item In either language, the \texttt{\$dbcIsExpZero} check (rational
  coefficient grouping by Boltzmann exponent) and the BFS queue management are
  the most straightforward components to port and should be implemented first
  as a proof-of-concept.
\end{enumerate}

%===============================================================
\section{Conclusions}
%===============================================================

SZ-DBC in its current form is a sound and largely correct tool for verifying
detailed balance in lattice MCMC algorithms.  The critical bugs from the first
audit have all been fixed.  The EBE-based exact check is rigorously correct, the
checker now always runs exactly (probabilistic SZ fallback requires explicit
opt-in), and the D4 check has been moved to opt-in status following the
performance analysis in Section~\ref{sec:d4perf}.  The result is a 1.53$\times$
speedup for \texttt{vmmc\_2d.wl} (332s $\to$ 217s) with no loss of DB correctness.

The remaining open issues are:
\begin{itemize}
\item A moderate correctness concern around EBE chamber reuse (unverified
  assumption, mathematically justified but not checked at runtime).
\item The \texttt{\$dbcFS} sentinel being silently treated as a DB violation
  (harmless for all standard algorithms tested).
\item A significant accuracy problem with the NormalDistribution discretisation
  (wrong transition probabilities for $\sigma \geq 1$, no warning printed).
\end{itemize}

For portability, the removal of FullSimplify from the critical path
substantially improves the prospects for an open-source reimplementation.  The
EBE inner loop requires only rational arithmetic, which is available in Python,
Julia, and C++.  The fundamental bottleneck for Python remains the lack of a
Block-equivalent for transparent RNG interception; Julia with Cassette.jl is the
recommended path for an open-source port that preserves the checker's
non-invasive design.

\end{document}
